\section{Introduction}\label{sec:introduction}
Let $p$ be a prime and let $R$ be a perfect $\F_p$-algebra or a nonzero
$p$-torsionfree perfectoid ring. Put
\[
X_R=\THH(R;\Z_p),\qquad
\pi_V^H X_R=[S^V,X_R]^H\quad(V\in RO(H)).
\]
We use the genuine commutative cyclotomic model of $\THH(HR)$
from \cite[Theorems~4.6--4.7]{NormTHH} and
\cite[Theorems~C--E]{BMCyclotomic}, with
\[
X_R=\holim_m\THH(HR)/p^m.
\]
The quotients are taken in spectra.
On geometric fixed points, a norm is a finite smash product over double cosets
\cite[Theorem~2.35]{NormTHH}. Thus norms preserve genuine mod-$p$ equivalences, and
\cite[Theorem~6.1]{BHNorm} implies the existence of norms on $X_R$.
On finite cyclic $p$-groups, isotropy separation and
\cite[Lemma~2.11(b)]{AN} show that geometric $C_p$-fixed points
commute with this uniformly connective tower.
Thus
\[
\varrho_p^*\Phi^{C_p}X_R\simeq X_R,
\qquad
\pi_k\Phi^{C_{p^i}}X_R=0
\quad(i\ge0,\ k<0).
\]
Thus $X_R$ admits a genuine cyclotomic restriction.


We can now state our main result.

\begin{theorem}\label{thm:main}
Let $p$ be a prime, and let $R$ be either a perfect
$\F_p$-algebra or a nonzero $p$-torsionfree perfectoid
ring. For every $G=C_{p^n}$, $n\ge0$, the $RO(G)$-graded
Tambara functor and genuine cyclotomic restriction of $X_R$ are given as follows.
\begin{enumerate}[label=\textup{(\roman*)}]
\item For nonzero $p$-torsionfree $R$, the homotopy groups $\pi_V^H X_R$ for every
$H\le G$ and $V\in RO(H)$ are given by
Proposition~\ref{thm:uniform}.
Products, restrictions, and transfers are given by
Proposition~\ref{prop:products-maps}, and Weyl actions
by Proposition~\ref{prop:weyl-action}.

Norms are given by
Propositions~\ref{prop:vertex-norms},
\ref{lem:period-repair-norm}, and
\ref{prop:edge-norms}.
The extension to arbitrary representation degrees
is given by Proposition~\ref{lem:ordered-sign-units}.
Genuine cyclotomic restriction is given by
Proposition~\ref{prop:cyclotomic-operations}.

\item For a perfect $\F_p$-algebra $R$, Theorem~\ref{thm:cp-result}
gives the graded Tambara functor by Euler localization of $\HZ$ and
extension of coefficients to $W(R)$, and gives genuine cyclotomic
restriction in every representation degree.
\end{enumerate}
\end{theorem}
The main strategy we use is Tate compression introduced in \cite[Theorem~4.7]{ZhangHZ}, so we write the homotopy groups as a kernel and a
cokernel.


Our calculation is related to several areas, including topological cyclic homology,
algebraic $K$-theory, and $p$-adic Hodge theory.

Our calculation applies to topological cyclic
homology. Write
\[
\TR^{n+1}(R;\Z_p)=X_R^{C_{p^n}}.
\]
For $n\ge1$, there are Frobenius, Verschiebung, and cyclotomic
restriction maps
\[
F,R:\TR^{n+1}(R;\Z_p)\longrightarrow
\TR^n(R;\Z_p),
\qquad
V:\TR^n(R;\Z_p)\longrightarrow
\TR^{n+1}(R;\Z_p).
\]
On homotopy groups, these maps are given by
$\res$, $R_{\mathrm{cyc}}$, and $\tr$, respectively.
For perfectoid rings, see \cite[Proposition~7.21]{MathewTR} for the
finite-level integer-graded $\TR$ rings and \cite[Theorem~3.7]{Andriopoulos}
for the formulas for $R$ and $F$.
For odd $p$ and $p$-torsionfree perfectoid $R$ containing a compatible
system of primitive $p$-power roots of unity, Davis and
Patchkoria~\cite[Section~8]{DavisPatchkoria} identify
$\pi_2\TR^n(R;\Z_p)$ with the Tate module $T_p(W_n\Omega_R^1)$.

For a perfect $\F_p$-algebra $R$, we have
\[
\pi_0\TR^{n+1}(R;\Z_p)\cong W_{n+1}(R),
\]
where $W_{n+1}(R)$ is the ring of Witt vectors of
length $n+1$~\cite[Theorem~3.3]{HM97}.
Our formulas extend this calculation to all real representation
degrees. In degree zero our norm is the Witt vector
norm of Angeltveit~\cite[Theorem~4.1]{WittNorm}.
For related degree-zero calculations using Witt vectors for Green
functors and Witt vectors with coefficients, see
\cite{BGHLWitt,DKNPWitt}.
For shifts by actual complex representations over perfect fields,
see \cite[Proposition~9.1]{HM97}.
For $R=\F_p$, the $\TR$-groups in complex and real circle representation
degrees were computed in \cite[Theorem~1.1]{GerhardtTR} and
\cite[Section~7]{AGTR}, respectively.

The spectrum $\THH(R;\Z_p)$ also appears in algebraic
$K$-theory. The cyclotomic trace~\cite{BHM} gives a map
\[
K(R)^\wedge_p\longrightarrow\operatorname{TC}(R;\Z_p).
\]
Topological cyclic homology is constructed from the
cyclotomic structure on $\THH(R;\Z_p)$
\cite{NS,BMCyclotomicOriginal}.
Our result gives the homotopy groups and the maps
in its finite cyclic fixed-point tower.
Hesselholt and Madsen~\cite{HMPolytopes} used representation
suspensions of $\THH(k)$ to compute $K_*(k[x]/(x^e),(x))$ for
a perfect field $k$ of positive characteristic.

For $p$-torsionfree perfectoid rings, our calculation
is also related to $p$-adic Hodge theory and
prismatic cohomology.
Let
\[
A=W(R^\flat),\qquad
I_0=\ker(\theta:A\to R),\qquad
I_j=\varphi^j(I_0).
\]
These ideals come from the perfect prism associated
to $R$~\cite[Theorem~3.10]{BSPrisms}. The Breuil--Kisin twist $B$ satisfies
\[
\pi_2\operatorname{TP}(R;\Z_p)\cong B,
\qquad
\pi_2\operatorname{TC}^-(R;\Z_p)\cong I_0B.
\]
For these isomorphisms, the canonical map, and linearized cyclotomic
Frobenius, see \cite[Proposition~6.5 and Remark~6.6]{BMS}
and Subsection~\ref{subsec:period-line}.
Proposition~\ref{thm:uniform} expresses the groups
$\pi_V^{C_{p^n}}X_R$ using $I_j$ and $B$.

For other historical works,
Hesselholt~\cite[Theorem~B]{HesCurves} described the integer-graded
homotopy groups of finite-level $\TR$ in terms of de Rham--Witt forms
and a polynomial generator in degree two for smooth algebras over a perfect field of characteristic $p$.
For perfectoid rings $R$, Sulyma studied the regular slice filtration
of $\THH(R;\Z_p)$~\cite{SulymaSlice} and the $RO(\mathbb T)$-graded
$\TF$ of $R$~\cite{SulymaTF}. He formulated the homotopy orbits to
$\TR$ spectral sequence as a spectral sequence of Mackey functors
and discussed the $RO(\mathbb T)$-graded Tambara structure of
$\THH(R;\Z_p)$~\cite[Sections~1 and~3.2]{SulymaTF}.

Let $A$ be a $\delta$-ring with a spherical lift $\mathbb S_A$.
For relatively quasiregular semiperfectoid $A$-algebras $S$,
Sulyma~\cite[Theorem~1.9]{SulymaNew} computed the representation-graded
groups of $\operatorname{TC}^-(S\mid\mathbb S_A)$ and
$\operatorname{TP}(S\mid\mathbb S_A)$.
For a polyrepresentation degree $\alpha$ in the contiguous range, he also computed
the form-degree $j$ contributions to finite-level relative $\TR$ in total degrees
$\alpha+j$ and $\alpha-1+j$ for $p$-completely
smooth algebras over $p$-torsionfree perfectoid rings and for smooth
algebras over perfect $\F_p$-algebras~\cite[Theorem~4.12]{SulymaNew}.
His norm formulas for de Rham--Witt forms apply in all form degrees
at odd primes and in even form degrees at $p=2$%
~\cite[Remark~1.20 and Theorem~5.7]{SulymaNew}.
He deferred a general treatment of Breuil--Kisin and prismatic norms
for the homotopy fixed points and Tate spectra of relative
$\THH$~\cite[Section~1.3.3]{SulymaNew}.
For the rings $R$ in Theorem~\ref{thm:main}, we determine the genuine
norms of $\THH(R;\Z_p)$ for every inclusion $H\le G=C_{p^n}$ and
every degree $V\in RO(H)$, including degrees with an odd coefficient
of the sign representation at $p=2$.

In characteristic $p$, Krause and Nikolaus~\cite[Proposition~2.1]{KrauseNikolaus}
identified the underlying spectrum of $\THH(k)$ with
$k\otimes_{\F_p}\THH(\F_p)$ for every perfect $\F_p$-algebra $k$.
At odd primes, Ayala, Mazel-Gee, and Rozenblyum~\cite[Corollary~10.8]{AMRDerived}
identified $\THH(\F_p)$ with $\Phi^{C_p}\HZ$ as genuine spectra on
the tower of cyclic $p$-groups. For the multiplicative comparison, see
Mao~\cite[Remark~7.6]{MaoDecomplete}.

Hu, Kriz, and Somberg~\cite[Theorem~2]{HKSMotivic} described the
equivariant motivic filtration on $\THH$ of polynomial $\F_p$-algebras.
Burghardt--Hu--Kriz--Kriz--Somberg \cite{Kriz}
gave chain-level models for equivariant $\THH$
in characteristic $p$.
Our characteristic-$p$ comparison extends their
module comparison to preserve norms.

Bhatt--Morrow--Scholze \cite[Theorem~1.6]{BMS} computed the
integer-graded homotopy groups of $\THH(-;\Z_p)$,
$\operatorname{TC}^-(-;\Z_p)$, and $\operatorname{TP}(-;\Z_p)$
for perfectoid rings.
For odd $p$ and $R=\mathcal O_{\mathbb C_p}$, the identification of
$\pi_0\operatorname{TC}^-(R;\Z_p)$ with $W(R^\flat)$ goes back
to Hesselholt~\cite{HesLocal}.
Mao~\cite{MaoThom} gave a Thom spectrum proof of B\"okstedt
periodicity for perfectoid rings.
Riggenbach \cite{Riggenbach} compared circle
homotopy fixed points with finite-subgroup
homotopy fixed points.
Bhatt--Lurie \cite{BhattLurie} developed the
divided conormal construction, which is used for the
Breuil--Kisin twist.

For real $\THH$ of perfectoid rings and for equivariant Witt complexes
in twisted $\THH$, see \cite{HornbostelPark} and \cite{BGKPY}, respectively.
Other work on cyclotomic spectra,
norms, and Tambara functors includes
\cite{AN,NS,NormTHH,BMCyclotomic,BMYRelative,
BHNorm,Ullman,Balderrama,HKSZOperations}.

Our conventions are as follows.
We work in complete universes. Fixed points, Tate constructions,
limits and smash products are derived.
For $V\in RO(G)$, $S^V$ is its representation sphere.
An ordered model of $V$ fixes the order of its irreducible summands. For $C_N$, fix
$\gamma_N=\exp(2\pi i/N)$; for $C_M\le C_N$, we have
$\gamma_M=\gamma_N^{N/M}$.

\begin{definition}\label{def:base-class}
Write $\mathbb T=S^1$ and let $C_{p^n}$ be its subgroup of
$p^n$th roots of unity. Let
\[
\varrho_p:\mathbb T\xrightarrow{\sim}\mathbb T/C_p
\]
be the inverse of $[z]\mapsto z^p$.
For $H\le G$, a degree $V\in RO(G)$ is \emph{natural} at $G/H$
if $\dim V^H\in2\mathbb Z$ and \emph{shifted} otherwise.
We omit the orbit when $H=G$.
\end{definition}

The $\Res,\Ind$ act on representations;
$\res,\tr,c_g,N$ denote restriction, transfer, conjugation and norm.
For degrees in $RO(G)$, $c_g$ also denotes the Weyl action.
We write $R_{\mathrm{cyc}}$ for cyclotomic restriction and
$F(-,-)$ for the function spectrum.

Write $\lambda_N(k)$ for the real plane of the complex character
of weight $k$, with ordered basis $(1,i)$. For $C_{p^n}$, put
\[
\lambda_j=\lambda_{p^n}(p^j)\quad(0\le j<n).
\]
At $p=2$, $\sigma_{2^m}$ is the sign representation of $C_{2^m}$
for $m\ge1$, modeled on $\mathbb R$ with basis $1$.
The identification $\lambda_{2^n}(2^{n-1})=2\sigma_{2^n}$ sends
the ordered basis $(1,i)$ to the two ordered real basis vectors.
For $n\ge1$, the \emph{selected degrees} of \cite{ZhangHZ} form
\[
RO_{\mathrm{sel}}(C_{p^n})=
\begin{cases}
\mathbb Z\{1,\lambda_0,\ldots,\lambda_{n-1}\},&p\ne2,\\
\mathbb Z\{1,\sigma_{2^n},\lambda_0,\ldots,\lambda_{n-2}\},&p=2,
\end{cases}
\qquad RO_{\mathrm{sel}}(e)=\mathbb Z.
\]
\begin{definition}\label{def:degree-coordinates}
For $V\in RO(C_{p^n})$, its degree coordinates are
$a_i=\frac12\dim_{\mathbb R}V^{C_{p^i}}$ $(0\le i\le n)$.
For natural degrees, $\mathbf a=(a_0,\ldots,a_n)$ belongs to
\[
D_n^{(0)}=\mathbb Z^{n+1},\qquad
D_n^{(1)}=(\mathbb Z+\tfrac12)^n\times\mathbb Z\ (n\ge1),\qquad
D_n=\begin{cases}
D_n^{(0)},&p\ne2\text{ or }n=0,\\
D_n^{(0)}\sqcup D_n^{(1)},&p=2,\ n\ge1.
\end{cases}
\]
Set $D_0^{(1)}=\varnothing$ and $\epsilon(\mathbf a)=\epsilon$ for $\mathbf a\in D_n^{(\epsilon)}$.
Define
\begin{equation*}
\rho_p(\mathbf a)=\begin{cases}
2a_n+\displaystyle\sum_{j=0}^{n-1}(a_j-a_{j+1})\lambda_j,
   &p\ne2,\ n\ge1,\\
2a_n+2(a_{n-1}-a_n)\sigma_{2^n}
 +\displaystyle\sum_{j=0}^{n-2}(a_j-a_{j+1})\lambda_j,
   &p=2,\ n\ge1,\\
2a_0,&n=0.
\end{cases}
\end{equation*}
\end{definition}
For $L=C_{p^m}$, write $\rho_L$ for $\rho_p$ formed over $L$.
For $p=2$, $n\ge1$ and a natural degree $U$,
$2(a_{n-1}-a_n)\equiv\epsilon(\mathbf a)\pmod2$.
For $\epsilon(\mathbf a)=0$, we lift $U$ to a virtual complex
$\mathbb T$-representation whose restriction to $C_{p^n}$ has underlying
real degree $U$. For selected degrees, we use weights $p^j$ and restrict
the same $\mathbb T$-representation to each subgroup.
The Euler class is $a_W\in\pi_{-W}^G\mathbb S$. For $H\le G$, set
$a_H^G=a_{\mathbb R[G/H]-1}$, and use the same notation for its
images in $X_R$.

We put $x^+=\max(x,0)$. A line over $A$ is an invertible $A$-module.
For a line $L$ and an integer $a$, set
\[
L^a=\begin{cases}L^{\otimes a},&a\ge0,\\
(L^\vee)^{\otimes(-a)},&a<0.
\end{cases}
\]
Frobenius is $\varphi$, with $\varphi^*L=A\otimes_{A,\varphi}L$.
A \emph{graded unit} of degree $W$ has a multiplicative inverse of degree $-W$.


The rest of this paper is organized as follows.
Section~\ref{sec:compression} applies Tate compression to $X_R$.
For nonzero $p$-torsionfree perfectoid rings, Section~\ref{sec:uniform}
computes the graded Tambara functor and cyclotomic restriction.

Appendix~\ref{app:spectral-foundations} computes the local homotopy
groups and comparison maps used in Section~\ref{sec:uniform}.
Appendix~\ref{app:operation-foundations} gives the remaining
proofs for norms.
We treat perfect $\F_p$-algebras in Appendix~\ref{app:characteristic-p},
using the calculations for $\HZ$ in \cite{ZhangHZ}.
